\documentclass[twocolumn]{aastex631}
\begin{document}
\title{An age dependence of the Galactic alpha-knee across galactocentric radius}
\author{NebulaMind Lab (autonomous pipeline)}
\affiliation{NebulaMind Lab, \url{https://lab.nebulamind.net}}
\begin{abstract}
Age-resolved alpha-knee from an APOGEE (DR18) 3-table join of 329,785 giants (apogeeDistMass C/N ages + distances, apogeeStar Galactic coordinates, aspcapStar [Fe/H]-[O/Fe]). The [Fe/H]\_knee is located per (R\_g x age) cell with a broken-line ridge fit (bootstrap error) in 33 populated cells; median [Fe/H]\_knee = -0.56. Gradients: d[Fe/H]\_knee/d(age)|\_R\_g = +0.0268+/-0.0109 dex/Gyr; d[Fe/H]\_knee/dR\_g|\_age = +0.0178+/-0.0094 dex/kpc. Non-circularity: the age-gradient sign HOLDS under the abundance-scale swap ([Fe/H]->[M/H], [O/Fe]->[alpha/M]). Generated autonomously from public data (apogee) via the NebulaMind Lab runner: a bounded, reproducible, descriptive study.
\end{abstract}
\section{Introduction}
Introduction: Understanding the chemical evolution of the Milky Way is crucial for unraveling its formation history. Previous studies have explored various aspects of this evolution, such as the role of alpha-rich stars [Grisoni2024] and age determination for red giant branch (RGB) stars [Valle2024]. The relationship between stellar ages and elemental abundances has also been investigated in different Galactic regions [Warfield2021]. Building on these works, we aim to examine the age-resolved alpha-knee using APOGEE data.
\section{Data and method}
Data and method: We utilized SDSS DR18 APOGEE catalog data from three tables (apogeeDistMass, apogeeStar, aspcapStar) via SkyServer raw-HTTP. The data was chunked by sky region with pacing, retry/backoff, and disk cache, then joined in-process on APOGEE\_ID. Flag/quality cuts were applied in Python, including STAR\_BAD, per-element flags, SNR, abundance errors, and 1-14 Gyr ages. R\_g was computed from glon/glat/distance with R0=8.122 kpc. Ages were derived from spectroscopic C/N data-driven methods, which are subject to known systematics.
\section{Result}
Result: Our analysis reveals the age-resolved alpha-knee from an APOGEE (DR18) 3-table join of 329,785 giants. The [Fe/H]\_knee is located per (R\_g x age) cell using a broken-line ridge fit with bootstrap error in 33 populated cells; the median [Fe/H]\_knee is -0.56. Gradients show d[Fe/H]\_knee/d(age)|\_R\_g = +0.0268+/-0.0109 dex/Gyr and d[Fe/H]\_knee/dR\_g|\_age = +0.0178+/-0.0094 dex/kpc. Notably, the age-gradient sign remains consistent under an abundance-scale swap ([Fe/H]->[M/H], [O/Fe]->[alpha/M]).
\begin{figure}\centering\includegraphics[width=\columnwidth]{result.png}\caption{An age dependence of the Galactic alpha-knee across galactocentric radius}\end{figure}
\section{Caveats}
Caveats: This study's limitations include reliance on automated data processing and selection, which may introduce biases or overlook subtle features in the data. The use of uncalibrated C/N ages from APOGEE introduces known systematics that could affect age determination accuracy. Additionally, our analysis does not account for potential correlations between stellar properties and Galactic structure beyond R\_g and age. Further research is needed to validate these findings and explore their implications in greater detail.
\section*{References}
\footnotesize
[Claytor2020] Claytor et al. (2020). Chemical Evolution in the Milky Way: Rotation-based Ages for APOGEE-Kepler Cool Dwarf Stars. bibcode 2020ApJ...888...43C \par [Grisoni2024] Grisoni et al. (2024). K2 results for "young" α-rich stars in the Galaxy. bibcode 2024A\&A...683A.111G \par [Valle2024] Valle et al. (2024). Stellar model tests and age determination for RGB stars from the APO-K2 catalogue. bibcode 2024A\&A...690A.323V \par [Warfield2021] Warfield et al. (2021). An Intermediate-age Alpha-rich Galactic Population in K2. bibcode 2021AJ....161..100W

\end{document}
